\documentclass[11pt,a4paper]{article}
\usepackage[T1]{fontenc}
\usepackage[utf8]{inputenc}
\usepackage[margin=18mm]{geometry}
\usepackage{amsmath,amssymb,booktabs,array,tabularx,hyperref,natbib,float}
\hypersetup{hidelinks}
\renewcommand{\thesection}{S\arabic{section}}
\renewcommand{\thetable}{S\arabic{table}}
\renewcommand{\arraystretch}{1.12}
\title{Supplementary Information\\[8pt]\large A Ground--Low-Altitude Mobility Manager for Service Reconfiguration under Disruptions}
\author{}
\date{}

% Revision v2: keep figures proportional and prevent float-page vertical spreading.
\usepackage{lmodern}
\usepackage{xurl}
\usepackage{placeins}
\newcolumntype{Y}{>{\raggedright\arraybackslash}X}
\emergencystretch=2em
\makeatletter
\setlength{\@fptop}{0pt}
\setlength{\@fpsep}{12pt}
\setlength{\@fpbot}{0pt plus 1fil}
\makeatother
\renewcommand{\topfraction}{0.95}
\renewcommand{\bottomfraction}{0.95}
\renewcommand{\textfraction}{0.05}
\renewcommand{\floatpagefraction}{0.80}
\setcounter{topnumber}{3}
\setcounter{totalnumber}{5}
\setlength{\textfloatsep}{12pt plus 2pt minus 2pt}
\setlength{\floatsep}{12pt plus 2pt minus 2pt}

\begin{document}
\maketitle
\section{Detailed experimental outcomes}

These tables retain the frozen results relocated from the main article. Table and figure captions in the main article identify the relevant populations. No experiments or statistical tests were rerun.

\begin{table}[H]
\centering
\small
\setlength{\tabcolsep}{3pt}
\caption{E1 service outcomes: mean [95\% interval].}
\label{tab:S1}
\begin{tabularx}{\textwidth}{@{}llrYYY@{}}
\toprule
Site & Policy & n & Completion (s) & Deadline violations (\%) & Damaged services/run \\
\midrule
A & B0 & 240 & 187.67 [183.63, 191.70] & 33.33 [27.67, 39.52] & 0.000 [0.000, 0.000] \\
A & B1 & 240 & 160.78 [154.28, 167.27] & 25.00 [19.94, 30.84] & 0.500 [0.436, 0.564] \\
A & B2 & 240 & 147.62 [142.15, 153.08] & 16.67 [12.48, 21.90] & 0.058 [0.028, 0.088] \\
A & B4b & 240 & 146.63 [141.41, 151.86] & 16.67 [12.48, 21.90] & 0.133 [0.090, 0.177] \\
B & B0 & 240 & 262.67 [249.67, 275.66] & 50.00 [43.72, 56.28] & 0.000 [0.000, 0.000] \\
B & B1 & 240 & 47.70 [46.62, 48.78] & 0.00 [0.00, 1.58] & 0.500 [0.436, 0.564] \\
B & B2 & 240 & 47.70 [46.62, 48.78] & 0.00 [0.00, 1.58] & 0.500 [0.436, 0.564] \\
B & B4b & 240 & 49.10 [47.24, 50.95] & 0.00 [0.00, 1.58] & 0.487 [0.424, 0.551] \\
C & B0 & 240 & 323.67 [318.37, 328.96] & 83.33 [78.10, 87.52] & 0.000 [0.000, 0.000] \\
C & B1 & 240 & 168.28 [166.83, 169.72] & 10.00 [6.81, 14.45] & 0.500 [0.436, 0.564] \\
C & B2 & 240 & 168.28 [166.83, 169.72] & 10.00 [6.81, 14.45] & 0.500 [0.436, 0.564] \\
C & B4b & 240 & 168.65 [167.07, 170.23] & 10.00 [6.81, 14.45] & 0.500 [0.436, 0.564] \\
\bottomrule
\end{tabularx}
\end{table}

Completion and damage use t intervals; deadline proportions use Wilson intervals. Each site--policy cell has 240 runs. These intervals characterize the fixed panel, not a population of cities.

\begin{table}[H]
\centering
\small
\setlength{\tabcolsep}{3pt}
\caption{E2 all-run transport outcomes and conditional path restoration.}
\label{tab:S2}
\begin{tabularx}{\textwidth}{@{}llrYYYY@{}}
\toprule
Site & Policy & All runs & Completion (s), mean [95\% CI] & Deadline violations (\%) & Recovered/affected & Recovery (\%), Wilson interval \\
\midrule
A & B0 & 320 & 182.00 [182.00, 182.00] & 100.00 & N/A & N/A \\
A & B1 & 320 & 162.00 [155.16, 168.84] & 37.50 & 200/200 & 100.00 [98.12, 100.00] \\
A & B2 & 320 & 162.00 [155.16, 168.84] & 37.50 & 200/200 & 100.00 [98.12, 100.00] \\
A & B4b & 320 & 166.53 [159.24, 173.81] & 39.06 & 191/191 & 100.00 [98.03, 100.00] \\
B & B0 & 320 & 241.00 [241.00, 241.00] & 100.00 & N/A & N/A \\
B & B1 & 320 & 153.88 [143.23, 164.52] & 37.50 & 200/200 & 100.00 [98.12, 100.00] \\
B & B2 & 320 & 153.88 [143.23, 164.52] & 37.50 & 200/200 & 100.00 [98.12, 100.00] \\
B & B4b & 320 & 154.05 [143.38, 164.72] & 37.50 & 199/199 & 100.00 [98.11, 100.00] \\
C & B0 & 320 & 310.00 [310.00, 310.00] & 100.00 & N/A & N/A \\
C & B1 & 320 & 278.00 [266.93, 289.07] & 62.50 & 200/200 & 100.00 [98.12, 100.00] \\
C & B2 & 320 & 278.00 [266.93, 289.07] & 62.50 & 200/200 & 100.00 [98.12, 100.00] \\
C & B4b & 320 & 278.92 [267.83, 290.02] & 62.81 & 201/201 & 100.00 [98.12, 100.00] \\
\bottomrule
\end{tabularx}
\end{table}

All-run denominators are 320 per site and policy. Recovery is conditional on actual primary-chain exposure; B0 is N/A. Recovery intervals are Wilson intervals.

\begin{table}[H]
\centering
\small
\setlength{\tabcolsep}{3pt}
\caption{Mean E3 system weighted loss by disturbance/task level.}
\label{tab:S3}
\begin{tabularx}{\textwidth}{@{}llrrrrr@{}}
\toprule
Site & Level & Runs/policy & B0 & B1 & B2 & B4b \\
\midrule
A & L1 & 40 & 240.000 & 0.000 & 0.000 & 6.000 \\
A & L2 & 80 & 240.000 & 240.000 & 240.000 & 240.000 \\
A & L3 & 80 & 240.000 & 0.000 & 0.000 & 1.688 \\
A & L4 & 40 & 240.000 & 120.000 & 120.000 & 123.375 \\
B & L1 & 40 & 240.000 & 0.000 & 0.000 & 0.000 \\
B & L2 & 80 & 240.000 & 120.000 & 120.000 & 120.000 \\
B & L3 & 80 & 375.000 & 135.000 & 135.000 & 135.000 \\
B & L4 & 40 & 375.000 & 255.000 & 255.000 & 255.000 \\
C & L1 & 40 & 240.000 & 240.000 & 240.000 & 240.000 \\
C & L2 & 80 & 240.000 & 240.000 & 240.000 & 240.000 \\
C & L3 & 80 & 375.000 & 135.000 & 135.000 & 135.000 \\
C & L4 & 40 & 375.000 & 375.000 & 375.000 & 375.000 \\
\bottomrule
\end{tabularx}
\end{table}

L1/L4 have 40 runs and L2/L3 have 80 per policy. Levels change disturbance structure and demand. SWL uses flat configured penalties and priority weights, rather than monetary costs.

\begin{table}[H]
\centering
\small
\setlength{\tabcolsep}{3pt}
\caption{E4A observation scheduling, completion and B4b input demand.}
\label{tab:S4}
\begin{tabularx}{\textwidth}{@{}rYYYYYYY@{}}
\toprule
Interval (s) & B0 completion (s) & B1 completion (s) & B2 completion (s) & B4b completion (s) & B4b late (\%) & B4b decisions/run & B4b prompt tokens/run \\
\midrule
10 & 192.0 & 140.0 & 140.0 & 140.0 & 0 & 62 & 172849.0 \\
30 & 212.0 & 150.0 & 150.0 & 150.0 & 0 & 22 & 62307.6 \\
60 & 253.0 & 180.0 & 180.0 & 180.0 & 0 & 12 & 35017.4 \\
120 & 253.0 & 240.0 & 240.0 & 240.0 & 100 & 7 & 21348.2 \\
300 & 482.0 & 521.0 & 482.0 & 482.0 & 100 & 4 & 11707.6 \\
\bottomrule
\end{tabularx}
\end{table}

Each policy--arm cell contains 20 runs. Decision counts cover the complete scheduled polling horizon, including post-resolution observations. Token means retain recorded backend outcomes.

\begin{table}[H]
\centering
\small
\setlength{\tabcolsep}{3pt}
\caption{E4B injected execution delay and mean completion duration.}
\label{tab:S5}
\begin{tabularx}{\textwidth}{@{}rrrrrr@{}}
\toprule
Delay (s) & B0 (s) & B1 (s) & B2 (s) & B4b (s) & B4b late (\%) \\
\midrule
0 & 182.0 & 120.0 & 120.0 & 120.0 & 0 \\
1 & 183.0 & 121.0 & 121.0 & 121.0 & 0 \\
5 & 187.0 & 125.0 & 125.0 & 125.0 & 0 \\
10 & 192.0 & 130.0 & 130.0 & 130.0 & 0 \\
20 & 202.0 & 140.0 & 140.0 & 140.0 & 0 \\
30 & 212.0 & 150.0 & 150.0 & 150.0 & 0 \\
60 & 302.0 & 341.0 & 302.0 & 302.0 & 100 \\
\bottomrule
\end{tabularx}
\end{table}

Each policy--arm cell contains 20 runs. Waiting times are controlled simulation interventions, not measured inference latency.

\begin{table}[H]
\centering
\small
\setlength{\tabcolsep}{3pt}
\caption{E4C B4b input burden with an unchanged core fleet.}
\label{tab:S6}
\begin{tabularx}{\textwidth}{@{}rrYYYY@{}}
\toprule
Total aircraft records & Core aircraft & Completion (s) & On time (\%) & Logged calls (20 runs) & Prompt tokens/logged call \\
\midrule
5 & 4 & 121.5 & 100 & 41 & 3982.1 \\
10 & 4 & 120.0 & 100 & 40 & 5920.6 \\
20 & 4 & 120.0 & 100 & 40 & 9588.7 \\
30 & 4 & 120.0 & 100 & 40 & 13256.7 \\
50 & 4 & 120.0 & 100 & 40 & 20593.5 \\
\bottomrule
\end{tabularx}
\end{table}

Each arm contains 20 B4b runs. Added records represent unavailable and non-commandable aircraft. Token means include retained error calls with zero recorded tokens; active fleet size stays four.

\begin{table}[H]
\centering
\small
\setlength{\tabcolsep}{3pt}
\caption{Aggregate E3 outcomes across the 12-scenario panel.}
\label{tab:S7}
\begin{tabularx}{\textwidth}{@{}llrYYY@{}}
\toprule
Site & Policy & n & Completion (s), mean [95\% CI] & Deadline violations (\%) & SWL, mean [95\% CI] \\
\midrule
A & B0 & 240 & 182.00 [182.00, 182.00] & 100.00 & 240.000 [240.000, 240.000] \\
A & B1 & 240 & 187.25 [176.00, 198.50] & 41.67 & 100.000 [84.923, 115.077] \\
A & B2 & 240 & 187.25 [176.00, 198.50] & 41.67 & 100.000 [84.923, 115.077] \\
A & B4b & 240 & 188.80 [177.39, 200.21] & 42.08 & 102.125 [86.923, 117.327] \\
B & B0 & 240 & 241.00 [241.00, 241.00] & 100.00 & 307.500 [298.899, 316.101] \\
B & B1 & 240 & 146.83 [133.44, 160.22] & 25.00 & 127.500 [113.165, 141.835] \\
B & B2 & 240 & 146.83 [133.44, 160.22] & 25.00 & 127.500 [113.165, 141.835] \\
B & B4b & 240 & 147.05 [133.67, 160.43] & 25.00 & 127.500 [113.165, 141.835] \\
C & B0 & 240 & 310.00 [310.00, 310.00] & 100.00 & 307.500 [298.899, 316.101] \\
C & B1 & 240 & 322.67 [306.37, 338.96] & 66.67 & 227.500 [217.182, 237.818] \\
C & B2 & 240 & 322.67 [306.37, 338.96] & 66.67 & 227.500 [217.182, 237.818] \\
C & B4b & 240 & 323.82 [307.38, 340.27] & 66.67 & 227.500 [217.182, 237.818] \\
\bottomrule
\end{tabularx}
\end{table}

Scenario--seed runs receive equal weight. L2/L3 therefore contribute twice as many runs as L1/L4; equal-level weighting would define a different summary.

\section{Interface and frozen implementation details}

\subsection{Principal notation}

\begin{table}[H]
\centering
\small
\setlength{\tabcolsep}{3pt}
\caption{Principal notation.}
\label{tab:S8}
\begin{tabularx}{\textwidth}{@{}lYY@{}}
\toprule
Symbol & Meaning & Unit \\
\midrule
$t,\tau(t),t_e,H$ & Decision time, observation timestamp, scheduled execution and horizon & s \\
$S_t,\widehat S_t$ & Current joint state and observed state & --- \\
$r_m,d_m,c_m$ & Mission release, deadline and completion & s \\
$\mathcal I_t,\mathcal C_{m,t}$ & Candidate information and legal subset & --- \\
$\widehat\eta_c,\delta_{\mathrm{exec}},\sigma_m$ & Remaining completion estimate, execution delay, estimated temporal margin & s \\
$a^{\mathrm{prop}},u,Q,\Gamma$ & Proposal, executed intervention, current-version indicator, current-state checker & --- \\
$J_t(c)$ & Frozen one-step heuristic score & Configured seconds-equivalent units \\
$A_i,R_i$ & Affected-chain and recovered-path indicators & Binary \\
$T_m,V_m,L_H$ & Completion duration, deadline violation, weighted service loss & s; binary; loss units \\
\bottomrule
\end{tabularx}
\end{table}

\subsection{Structured action vocabulary}

\begin{table}[H]
\centering
\small
\setlength{\tabcolsep}{3pt}
\caption{Structured action vocabulary.}
\label{tab:S9}
\begin{tabularx}{\textwidth}{@{}lY@{}}
\toprule
Action & Meaning \\
\midrule
DISPATCH & Assign an available compatible aircraft \\
REASSIGN & Transfer an eligible occupied aircraft to another mission \\
REROUTE & Change an ongoing operation's route \\
DIVERT & Change an air operation's destination or continuation \\
DELAY & Postpone an operation \\
CANCEL & Cancel a specified operation \\
RESERVE & Reserve a resource for a mission \\
RETURN & Request return by an eligible aircraft \\
LAND & Request landing under modeled constraints \\
GROUND\_FALLBACK & Establish ground service for an eligible mission \\
NO\_ACTION & Make no operational change at the decision epoch \\
ESCALATE & Report a need beyond the current action choice \\
\bottomrule
\end{tabularx}
\end{table}

The vocabulary is larger than the set exercised by any single scenario. Supporting an action type in the contract does not imply that the experiments independently validate every possible use of that action. The candidate interface retains legal flags and rejection reasons, while the checker applies action-specific conditions.

\subsection{Frozen policy details}

The prompt template is \nolinkurl{manager_v2.txt}, SHA-256 \nolinkurl{cf3a550761718d4b5ba53cd004980d795cc8e0ac9708f4284fe0f22be94f279c}. B4a and B4b use the same template, with only the candidate section withheld in B4a. B2's weight file has SHA-256 \nolinkurl{a6e880f3c8e7398bcc2c25568141bd141be24d04eab7a8d0989cc72cacc673f6}. Its preemption penalty depends on the incumbent mission's priority. These values were not retuned for this manuscript. The backend identifier reported in Section 4.3 is preserved as provenance, not interpreted as an independently verified model-release label.

\section{Statistical and outcome conventions}

\subsection{Conditional populations}

Recovery outcomes apply only after the primary service chain is affected. At Site C in E2, B4b has 201 affected runs, but the B4b--B2 recovery-time comparison contains 200 jointly affected pairs. All-run completion and deadline comparisons contain 320 pairs per site. E4A OBS300 is unexposed at the time of failure and has no conditional recovery denominator. Its late completions remain in the all-run outcome analysis.

\subsection{Comparison families and missing fields}

The E2 B4b--B2 family contains completion time, recovery time, failure-to-replan latency, existing-service damage, candidate-set reduction, deadline violation, recovery success, recovered-and-completed status, ground fallback, failure-induced ground fallback, necessary-ground-fallback correctness and air intervention. The candidate-count endpoint mixes occupancy and availability changes and is not used as a clean causal estimate of failure-induced candidate loss. Undefined or unavailable tests remain conservatively represented in the recorded adjustment family without turning missing effects into estimated zero effects.

E3 compound comparisons cover L2--L4, B1 and B2 as comparators, and SWL and CRITICAL+HIGH loss as endpoints. Action-level priority-consistency and resource-competition fields are unimplemented and are not analyzed as zero violations or perfect correctness. The manuscript uses observed priority-weighted outcomes instead. E4 uses separate families for each policy, metric and subexperiment, retaining the original anchors.

\section{Additional timing and implementation records}

The zero-delay E4B anchor matches the corresponding E2 transport outcomes for all 80 policy--seed runs. B0/B1/B2 action logs match in 20/20 cases and B4b in 19/20. The remaining B4b log contains an additional retained transport error. Backend wall time does not advance physical simulation time.

Switching resources or modes is distinct from returning to a previously used choice. A necessary replacement aircraft alone does not establish decision oscillation. In E4A, prompt demand falls from approximately 172,849 to 11,708 tokens per run across OBS10--OBS300, under the fixed polling horizon.

The analysis environment recorded in the final release is distinct from individual run configurations. Its source snapshot and hashes support traceability without retroactively proving identical historical environments.

\section{Figure production and source records}

Figure 1 combines separate, text-free AI-generated illustrations with editable PowerPoint text, routes, containers and markers. Its map, scale and four findings are schematic. The visual reference was supplied by the authors; generated illustrations do not depict observations or operational deployments at the study sites. Images were produced with the built-in imagegen tool and composed into PowerPoint. All scientific labels and claims were checked against the frozen reports.

Figures 2--8 were restyled from the existing editable sources. Road geometry and plotted numerical series remain unchanged. Quantitative charts retain embedded workbooks; their series are rounded to eight decimal places for workbook portability, while manuscript tables retain the frozen released values. Policy colors identify policies; route colors identify transport modes. Reproducible build sources and verification records accompany the local delivery.

\FloatBarrier
\end{document}
